\section*{LỜI MỞ ĐẦU}
\phantomsection\addcontentsline{toc}{section}{\numberline {}LỜI MỞ ĐẦU}
Trong bối cảnh chuyển đổi số đang diễn ra mạnh mẽ, công nghệ thông tin đã trở thành nền tảng quan trọng trong hầu hết các lĩnh vực của đời sống xã hội. Một trong những ứng dụng tiêu biểu là bầu cử điện tử, nơi các quy trình đăng ký, bỏ phiếu và kiểm phiếu được thực hiện qua môi trường số, nâng cao tính thuận tiện và khả năng tiếp cận của cử tri. Tuy nhiên, bầu cử có yêu cầu rất cao về tính minh bạch, công bằng, toàn vẹn dữ liệu và bảo mật thông tin cá nhân. Việc ứng dụng blockchain vào bầu cử điện tử mở ra hướng tiếp cận phù hợp nhờ đặc tính phân tán, khó sửa đổi, có thể kiểm chứng công khai và giảm sự phụ thuộc vào một trung tâm quản lý duy nhất.

Các hệ thống bầu cử điện tử truyền thống tồn tại nhiều hạn chế. Dữ liệu phiếu bầu được lưu trữ tập trung nên có nguy cơ bị tấn công hoặc chỉnh sửa nếu máy chủ quản lý bị xâm nhập. Quy trình kiểm phiếu phụ thuộc vào đơn vị vận hành, khiến cử tri khó tự mình xác minh tính đúng đắn của kết quả. Ngoài ra, hệ thống cũ còn gặp thách thức trong việc đảm bảo đồng thời hai yêu cầu tưởng như mâu thuẫn: xác thực đúng tư cách cử tri và giữ bí mật tuyệt đối lựa chọn của họ, dẫn đến các rủi ro như bỏ phiếu trùng lặp, giả mạo danh tính hoặc lộ thông tin cá nhân nhạy cảm.

Từ những vấn đề trên, em đề xuất đề tài \textbf{Xây dựng hệ thống bỏ phiếu điện tử sử dụng công nghệ chuỗi khối}. Đề tài hướng đến xây dựng mô hình bầu cử điện tử có tính an toàn, minh bạch và khả năng kiểm chứng cao, dựa trên hai công nghệ cốt lõi: \textbf{Hyperledger Fabric Blockchain} và \textbf{chữ ký số tập thể mù dựa trên lược đồ EC-Schnorr}. Blockchain cho phép ghi nhận phiếu bầu dưới dạng các giao dịch bất biến, đảm bảo toàn vẹn dữ liệu và hạn chế nguy cơ chỉnh sửa trái phép. Đồng thời, cơ chế chữ ký số tập thể mù xác thực quyền bỏ phiếu mà không làm lộ nội dung lá phiếu, bảo vệ tính riêng tư và ẩn danh của cử tri. Sự kết hợp này góp phần nâng cao khả năng chống gian lận, tăng cường tính minh bạch và củng cố niềm tin vào kết quả bầu cử.

Mục tiêu của đề tài là xây dựng hệ thống hỗ trợ các chức năng cốt lõi gồm đăng ký cử tri, xác thực tư cách tham gia, bỏ phiếu trực tuyến, ghi nhận phiếu bầu lên blockchain, kiểm phiếu tự động và xác minh phiếu bầu sau bầu cử. Đề tài hướng đến một giải pháp vừa đáp ứng yêu cầu kỹ thuật về bảo mật, vừa phù hợp với nhu cầu thực tiễn trong tổ chức bầu cử hiện đại.
\newpage
\clearpage
